\section{Data and evidence stages}
\label{sec:data}
\begin{table*}[t]
\centering\small
\caption{Six scored sources. An origin is a forecasting time, not an independent
organizational replication. New-period experiments are selected after development
pilots; their models are fixed before target preparation and scoring.}
\label{tab:sources}
\begin{tabular}{llllr}
\toprule
Source & Current-cohort target & Observed relations & Evidence stage & Origins\\
\midrule
Borg 2011 & First pending/running state exit & Job, machine & Original frozen period & 28\\
BPI 2017 & First qualifying application milestone & Last handler & Original frozen period & 61\\
BPI 2012 & First approval/decline/cancellation & Last handler & First confirmation & 28\\
Azure V2 & VM deletion & Deployment, subscription & New-period confirmation & 28\\
UCI incidents & First recorded resolution/closure & Handler, team & New-period confirmation & 28\\
BPI 2014 & First Closed activity after Open & Team & New-period confirmation & 28\\
\bottomrule
\end{tabular}
\end{table*}
The sources span two compute providers and several workflow/incident releases
(Table~\ref{tab:sources}). Their horizons and target semantics differ; scores
are compared within source, not pooled into one cross-domain leaderboard.
All features and memberships use recorded events at or before the origin.
Entire outcome windows must mature before their partition ends. Historical
extracts do not establish complete live populations or ingestion availability.

The initial Borg/BPI 2017 study froze 264 assets. Its outcomes motivated
retrospective direct-count, calibration, and joint-risk controls. BPI 2012 then
froze 82 assets before February scoring. The new extension freezes 165 assets
for Azure days 22--28.75, UCI May 1--28, 2016, and BPI 2014 January 1--28,
2014. The latter adds 84 origins without retuning on their scores. These are
internal holdout freezes, not externally preregistered studies.

New-source model fitting, marginal calibration, dependence fitting, development
validation, and confirmation occupy distinct chronological windows. Three-seed
neural/boosted-tree CDF ensembles are selected by earlier calibration Brier;
copulas and direct counts use the same dependence window. Controls include
feature-matched temporal HGB, node-only and global-set networks, typed message
passing, temporal attention, and relation-specific/shuffled graphs. Gaussian
and Student-$t$ couplings retain the selected marginal probabilities. Complete
windows, feature exclusions, model budgets, and source-selection disclosures
appear in Appendix~\ref{app:design}.

Azure retains all released VMs in a deterministic 5\% sample of subscriptions,
not a prefix of arbitrary VM rows. Google 2019 is audited on one shard only;
incomplete trajectories prevent a valid predictive benchmark. BPI change
snapshots lack reliable attribute-version histories. Citi Bike and two airline
geographies are exploratory pilots, not additional confirmation sources.
